\begin{textsample}{NEG Dim 8 – human – Score -2.00 – rj/rj\_ef\_ciencias\_da\_natureza\_1\_p3.txt}  \label{ex:f8_neg_006}
No trabalho de construção do DOC-RJ \textbf{elaborado} de maneira alinhada à Base Nacional Comum Curricular privilegia -se o principal sujeito da ação educativa : o estudante.
Para isso, o professor precisa oportunizar situações de protagonismo, assumindo o seu papel de mediador e apropriando -se de metodologias de aprendizagem ativa.
A metodologia ativa é uma concepção educativa que estimula processos construtivosdeação-reflexão-ação(FREIRE,2006), em que o estudante \textbf{participa} ativamente de todo processo.
Por \textbf{exemplo}, a Aprendizagem por Resolução de Problemas- ARP ( Metodologia da Problematização ).
Na ARP os estudantes são estimulados a observar em sua volta e perceberem os problemas que podem ser estudados e solucionados.
De acordo com a BNCC ( 2017 ), o processo investigativo deve ser entendido como elemento chave na formação dos estudantes.
Para tanto, o aprendizado deve ir além do passo a passo e do cumprimento das etapas pré-definidas do método científico. É indispensável o estímulo à observação, experimentação e investigação além da simples execução de um roteiro laboratorial, a fim de formar estudantes questionadores e divulgadores do conhecimento científico.

% matched lemmas: elaborar, exemplo, participar
\end{textsample}
